% Figure: absorption budget for the auditorium case study. Stacked
% contribution of each surface class to the total sabins at 500 Hz, with the
% occupied-audience contribution dominating.
\begin{figure}[htbp]
\centering
\begin{tikzpicture}
\begin{axis}[
  width=12cm, height=7.5cm,
  ybar,
  bar width=16pt,
  xlabel={surface class},
  ylabel={absorption (m$^2$-sabin at 500\,Hz)},
  ymin=0, ymax=450,
  axis lines=left,
  symbolic x coords={audience,seats,ceiling,sidewalls,floor,rear,stage},
  xtick=data,
  x tick label style={rotate=35, anchor=east, font=\scriptsize},
  tick label style={font=\scriptsize},
  label style={font=\small},
  nodes near coords,
  nodes near coords style={font=\tiny, /pgf/number format/fixed, /pgf/number format/precision=0},
  every node near coord/.append style={anchor=south},
  grid=major, grid style={enggray!20},
]
\addplot[fill=engblue!70, draw=engblue] coordinates {
  (audience,382.5) (seats,90) (ceiling,40) (sidewalls,60)
  (floor,18) (rear,120) (stage,55)
};
\end{axis}
\end{tikzpicture}
\caption[Absorption budget of the auditorium by surface class]{Absorption
budget of the case-study auditorium at $500\,\si{\hertz}$, broken out by
surface class. The occupied audience dominates the total sabins, which is
the controlling fact of large-hall acoustics: people are the largest
absorber, so a hall that sounds correct when full sounds too live when empty
unless the unoccupied seats are made to absorb almost as much as a seated
listener. The rear-wall treatment is the next-largest term, sized to kill
the long-path echo back to the stage. The hard reflecting surfaces, the
ceiling and the side walls near the stage, contribute little absorption by
design, because their job is early-reflection support rather than decay
control.}
\label{fig:vol03ch07:absorption-budget}
\end{figure}
